\documentclass[12pt,a4paper]{article}
\usepackage{fontspec}
\usepackage[russian]{babel}
\usepackage{amsmath,amssymb,amsfonts,amsthm}
\usepackage{geometry}
\usepackage{enumitem}
\usepackage{microtype}

\geometry{margin=2cm}
\defaultfontfeatures{Ligatures=TeX}
\setmainfont{Times New Roman}
\setlength{\parindent}{0pt}
\setlength{\parskip}{0.45em}
\setlist[enumerate]{itemsep=0.55em, topsep=0.7em, parsep=0.25em}
\setlength{\abovedisplayskip}{0.8em}
\setlength{\belowdisplayskip}{0.8em}

\begin{document}

\section*{Задание 3. Поверхности}

\textit{Ниже нормаль в каждой задаче берется как $\mathbf m = \dfrac{\mathbf r_u \times \mathbf r_v}{\|\mathbf r_u \times \mathbf r_v\|}$. При смене ориентации нормали меняются знак второй фундаментальной формы, главных кривизн, средней кривизны и нормальной кривизны, но не гауссова кривизна.}

\subsection*{Формулы, которые используются}

\begin{enumerate}
    \item \textbf{Поверхность вращения.}
    Если кривая профиля имеет вид $\gamma(u) = (x(u),0,z(u))$, то
    \[
    \mathbf r(u,v) = \bigl(x(u)\cos v,\ x(u)\sin v,\ z(u)\bigr).
    \]

    \item \textbf{Базис касательного пространства.}
    \[
    \mathbf r_u = \frac{\partial \mathbf r}{\partial u}, \qquad
    \mathbf r_v = \frac{\partial \mathbf r}{\partial v}.
    \]

    \item \textbf{Единичная нормаль.}
    \[
    \mathbf m = \frac{\mathbf r_u \times \mathbf r_v}{\|\mathbf r_u \times \mathbf r_v\|}.
    \]

    \item \textbf{Касательная плоскость и нормальная прямая.}
    В точке $P$:
    \[
    \mathbf m(P)\cdot\bigl((X,Y,Z)-P\bigr)=0, \qquad (X,Y,Z)=P+\lambda\,\mathbf m(P).
    \]

    \item \textbf{Первая фундаментальная форма.}
    \[
    I = E\,du^2 + 2F\,du\,dv + G\,dv^2,
    \]
    где
    \[
    E=\langle \mathbf r_u,\mathbf r_u\rangle,\qquad
    F=\langle \mathbf r_u,\mathbf r_v\rangle,\qquad
    G=\langle \mathbf r_v,\mathbf r_v\rangle.
    \]

    \item \textbf{Вторая фундаментальная форма.}
    \[
    II = L\,du^2 + 2M\,du\,dv + N\,dv^2,
    \]
    где
    \[
    L=\langle \mathbf r_{uu},\mathbf m\rangle,\qquad
    M=\langle \mathbf r_{uv},\mathbf m\rangle,\qquad
    N=\langle \mathbf r_{vv},\mathbf m\rangle.
    \]

    \item \textbf{Гауссова кривизна.}
    \[
    K = \frac{LN-M^2}{EG-F^2}.
    \]

    \item \textbf{Нормальная кривизна.}
    Для касательного вектора $a=a^1\mathbf r_u+a^2\mathbf r_v$:
    \[
    k_n(a)=\frac{II(a)}{I(a)}
    =\frac{L(a^1)^2+2Ma^1a^2+N(a^2)^2}{E(a^1)^2+2Fa^1a^2+G(a^2)^2}.
    \]
    В частности, для $a=3\mathbf r_u-2\mathbf r_v$:
    \[
    k_n=\frac{9L-12M+4N}{9E-12F+4G}.
    \]

    \item \textbf{Главные кривизны и средняя кривизна.}
    \[
    \det(B-\lambda G)=0, \qquad H=\frac{k_1+k_2}{2}, \qquad K=k_1k_2.
    \]

    \item \textbf{График функции $z=f(x,y)$.}
    Для поверхности
    \[
    \mathbf r(x,y)=(x,y,f(x,y))
    \]
    используются формулы
    \[
    \mathbf m=\frac{(-f_x,-f_y,1)}{\sqrt{1+f_x^2+f_y^2}},\qquad
    E=1+f_x^2,\quad F=f_xf_y,\quad G=1+f_y^2,
    \]
    \[
    L=\frac{f_{xx}}{\sqrt{1+f_x^2+f_y^2}},\quad
    M=\frac{f_{xy}}{\sqrt{1+f_x^2+f_y^2}},\quad
    N=\frac{f_{yy}}{\sqrt{1+f_x^2+f_y^2}},\quad
    K=\frac{f_{xx}f_{yy}-f_{xy}^2}{(1+f_x^2+f_y^2)^2}.
    \]
\end{enumerate}

\subsection*{а) Поверхность вращения кривой $\gamma(t) = (x(t), 0, z(t))$ вокруг оси $z$}

Обозначим параметр профиля через $u$, а угол вращения через $v$. Тогда
\[
\mathbf r(u,v) = \bigl(x(u)\cos v,\ x(u)\sin v,\ z(u)\bigr).
\]

\textbf{1) Параметризация, базис касательного пространства, нормаль}
\[
\mathbf r_u = \bigl(x'(u)\cos v,\ x'(u)\sin v,\ z'(u)\bigr), \qquad
\mathbf r_v = \bigl(-x(u)\sin v,\ x(u)\cos v,\ 0\bigr).
\]

При $x(u)>0$
\[
\mathbf m = \frac{\mathbf r_u \times \mathbf r_v}{\|\mathbf r_u \times \mathbf r_v\|}
= \frac{\bigl(-z'(u)\cos v,\ -z'(u)\sin v,\ x'(u)\bigr)}{\sqrt{(x'(u))^2 + (z'(u))^2}}.
\]

\textbf{2) Касательная плоскость и нормальная прямая}

Пусть $P=\mathbf r(u_0,v_0)=(x_0\cos v_0,x_0\sin v_0,z_0)$, где $x_0=x(u_0)$, $z_0=z(u_0)$.
Тогда касательная плоскость:
\[
\mathbf m(u_0,v_0)\cdot \bigl((X,Y,Z)-P\bigr)=0.
\]
В развернутом виде:
\[
-z'(u_0)\cos v_0\,(X-x_0\cos v_0) - z'(u_0)\sin v_0\,(Y-x_0\sin v_0)
+ x'(u_0)\,(Z-z_0)=0.
\]

Нормальная прямая:
\[
(X,Y,Z)=P+\lambda\,\bigl(-z'(u_0)\cos v_0,\ -z'(u_0)\sin v_0,\ x'(u_0)\bigr).
\]

\textbf{3) Первая и вторая фундаментальные формы, гауссова кривизна}

\[
E=(x')^2+(z')^2,\qquad F=0,\qquad G=x^2.
\]
\[
L=\frac{x'z''-x''z'}{\sqrt{(x')^2+(z')^2}},\qquad M=0,\qquad
N=\frac{xz'}{\sqrt{(x')^2+(z')^2}}.
\]

Следовательно,
\[
K=\frac{LN-M^2}{EG-F^2}
=\frac{z'(x'z''-x''z')}{x\bigl((x')^2+(z')^2\bigr)^2}.
\]

\textbf{4) Нормальная кривизна в направлении $a=3\mathbf r_u-2\mathbf r_v$}

\[
k_n(a)=\frac{\mathrm{II}(a)}{\mathrm{I}(a)}
=\frac{9L-12M+4N}{9E-12F+4G}
=\frac{9L+4N}{9E+4G}.
\]
То есть
\[
k_n(a)=
\frac{9(x'z''-x''z')+4xz'}{\sqrt{(x')^2+(z')^2}\,\bigl(9((x')^2+(z')^2)+4x^2\bigr)}.
\]

\textbf{5) Главные кривизны и средняя кривизна}

Так как $F=M=0$, координатные линии являются главными направлениями. Поэтому
\[
k_1=\frac{L}{E}=\frac{x'z''-x''z'}{\bigl((x')^2+(z')^2\bigr)^{3/2}},\qquad
k_2=\frac{N}{G}=\frac{z'}{x\sqrt{(x')^2+(z')^2}}.
\]

\[
H=\frac{k_1+k_2}{2},\qquad K=k_1k_2.
\]

\subsection*{b) Сфера: $\gamma(t) = (R\cos t, 0, R\sin t)$, $u_0=\frac{\pi}{4}$, $v_0=\frac{\pi}{4}$}

Здесь
\[
\mathbf r(u,v)=\bigl(R\cos u\cos v,\ R\cos u\sin v,\ R\sin u\bigr).
\]
Точка:
\[
P=\mathbf r\!\left(\frac{\pi}{4},\frac{\pi}{4}\right)=
\left(\frac R2,\ \frac R2,\ \frac{R\sqrt2}{2}\right).
\]

\textbf{1) Базис касательного пространства и нормаль}

\[
\mathbf r_u=\bigl(-R\sin u\cos v,\ -R\sin u\sin v,\ R\cos u\bigr),\qquad
\mathbf r_v=\bigl(-R\cos u\sin v,\ R\cos u\cos v,\ 0\bigr).
\]
В точке $P$:
\[
\mathbf r_u(P)=\left(-\frac R2,-\frac R2,\frac{R\sqrt2}{2}\right),\qquad
\mathbf r_v(P)=\left(-\frac R2,\frac R2,0\right).
\]

\[
\mathbf m(u,v)=\frac{\mathbf r_u\times \mathbf r_v}{\|\mathbf r_u\times \mathbf r_v\|}
=-\bigl(\cos u\cos v,\ \cos u\sin v,\ \sin u\bigr),
\]
поэтому
\[
\mathbf m(P)=\left(-\frac12,-\frac12,-\frac1{\sqrt2}\right).
\]

\textbf{2) Касательная плоскость и нормальная прямая}

\[
-\frac12\left(x-\frac R2\right)-\frac12\left(y-\frac R2\right)-\frac1{\sqrt2}\left(z-\frac{R\sqrt2}{2}\right)=0.
\]
После упрощения:
\[
x+y+\sqrt2\,z=2R.
\]

\[
(x,y,z)=\left(\frac R2,\frac R2,\frac{R\sqrt2}{2}\right)+
\lambda\left(-\frac12,-\frac12,-\frac1{\sqrt2}\right).
\]

\textbf{3) Первая и вторая фундаментальные формы, гауссова кривизна}

\[
E=R^2,\qquad F=0,\qquad G=R^2\cos^2 u,
\]
в точке $u_0=\pi/4$:
\[
E=R^2,\qquad F=0,\qquad G=\frac{R^2}{2}.
\]

\[
L=R,\qquad M=0,\qquad N=R\cos^2 u,
\]
в точке:
\[
L=R,\qquad M=0,\qquad N=\frac R2.
\]

\[
K=\frac{LN-M^2}{EG-F^2}=\frac1{R^2}.
\]

\textbf{4) Нормальная кривизна в направлении $a=3\mathbf r_u-2\mathbf r_v$}

\[
k_n(a)=\frac{9L-12M+4N}{9E-12F+4G}
=\frac{9R+4R\cos^2 u}{9R^2+4R^2\cos^2 u}
=\frac1R.
\]
В частности, в точке $\left(\frac{\pi}{4},\frac{\pi}{4}\right)$:
\[
k_n(a)=\frac1R.
\]

\textbf{5) Главные кривизны и средняя кривизна}

\[
k_1=\frac{L}{E}=\frac1R,\qquad k_2=\frac{N}{G}=\frac1R,\qquad
H=\frac{k_1+k_2}{2}=\frac1R.
\]

\subsection*{c) Тор: $\gamma(t) = (R+\rho\cos t, 0, \rho\sin t)$, $u_0=\frac{\pi}{6}$, $v_0=\frac{\pi}{3}$}

Чтобы не путать параметризацию и радиус образующей окружности, обозначим его через $\rho$; это тот же самый $r$ из условия.

\[
\mathbf r(u,v)=\bigl((R+\rho\cos u)\cos v,\ (R+\rho\cos u)\sin v,\ \rho\sin u\bigr).
\]
Положим
\[
A=R+\frac{\sqrt3}{2}\rho.
\]
Тогда
\[
P=\mathbf r\!\left(\frac{\pi}{6},\frac{\pi}{3}\right)=
\left(\frac A2,\ \frac{\sqrt3\,A}{2},\ \frac{\rho}{2}\right).
\]

\textbf{1) Базис касательного пространства и нормаль}

\begin{align*}
\mathbf r_u&=\bigl(-\rho\sin u\cos v,\ -\rho\sin u\sin v,\ \rho\cos u\bigr),\\
\mathbf r_v&=\bigl(-(R+\rho\cos u)\sin v,\ (R+\rho\cos u)\cos v,\ 0\bigr).
\end{align*}
В точке $P$:
\[
\mathbf r_u(P)=\left(-\frac{\rho}{4},\ -\frac{\sqrt3\,\rho}{4},\ \frac{\sqrt3\,\rho}{2}\right),\qquad
\mathbf r_v(P)=\left(-\frac{\sqrt3\,A}{2},\ \frac A2,\ 0\right).
\]

\[
\mathbf m(u,v)=\frac{\mathbf r_u\times \mathbf r_v}{\|\mathbf r_u\times \mathbf r_v\|}
=-\bigl(\cos u\cos v,\ \cos u\sin v,\ \sin u\bigr),
\]
поэтому
\[
\mathbf m(P)=\left(-\frac{\sqrt3}{4},\ -\frac34,\ -\frac12\right).
\]

\textbf{2) Касательная плоскость и нормальная прямая}

\[
-\frac{\sqrt3}{4}(x-x_0)-\frac34(y-y_0)-\frac12(z-z_0)=0,
\]
где
\[
x_0=\frac A2,\qquad y_0=\frac{\sqrt3\,A}{2},\qquad z_0=\frac{\rho}{2}.
\]
После упрощения:
\[
\sqrt3\,x+3y+2z=2\sqrt3\,R+4\rho.
\]

\[
(x,y,z)=\left(\frac A2,\ \frac{\sqrt3\,A}{2},\ \frac{\rho}{2}\right)
+\lambda\left(-\frac{\sqrt3}{4},\ -\frac34,\ -\frac12\right).
\]

\textbf{3) Первая и вторая фундаментальные формы, гауссова кривизна}

\[
E=\rho^2,\qquad F=0,\qquad G=(R+\rho\cos u)^2,
\]
а в точке $u_0=\pi/6$:
\[
E=\rho^2,\qquad F=0,\qquad G=A^2.
\]

\[
L=\rho,\qquad M=0,\qquad N=(R+\rho\cos u)\cos u,
\]
в точке:
\[
L=\rho,\qquad M=0,\qquad N=\frac{\sqrt3}{2}A.
\]

\[
K=\frac{LN-M^2}{EG-F^2}
=\frac{\cos u}{\rho(R+\rho\cos u)}.
\]
В точке $\left(\frac{\pi}{6},\frac{\pi}{3}\right)$:
\[
K=\frac{\sqrt3}{2\rho\left(R+\frac{\sqrt3}{2}\rho\right)}
=\frac{\sqrt3}{2\rho R+\sqrt3\,\rho^2}.
\]

\textbf{4) Нормальная кривизна в направлении $a=3\mathbf r_u-2\mathbf r_v$}

\[
k_n(a)=\frac{9L-12M+4N}{9E-12F+4G}
=\frac{9\rho+4(R+\rho\cos u)\cos u}{9\rho^2+4(R+\rho\cos u)^2}.
\]
В точке $\left(\frac{\pi}{6},\frac{\pi}{3}\right)$:
\[
k_n(a)=\frac{12\rho+2\sqrt3\,R}{9\rho^2+4\left(R+\frac{\sqrt3}{2}\rho\right)^2}.
\]

\textbf{5) Главные кривизны и средняя кривизна}

Так как $F=M=0$, главные направления совпадают с меридианами и параллелями:
\[
k_1=\frac{L}{E}=\frac1{\rho},\qquad
k_2=\frac{N}{G}=\frac{\cos u}{R+\rho\cos u}.
\]

В точке $\left(\frac{\pi}{6},\frac{\pi}{3}\right)$:
\[
k_1=\frac1{\rho},\qquad
k_2=\frac{\sqrt3}{2R+\sqrt3\,\rho}.
\]

\[
H=\frac{k_1+k_2}{2}.
\]
В точке:
\[
H=\frac12\left(\frac1{\rho}+\frac{\sqrt3}{2R+\sqrt3\,\rho}\right).
\]

\subsection*{d) График функции $z=f(x,y)$}

\[
\mathbf r(x,y)=\bigl(x,\ y,\ f(x,y)\bigr).
\]

\textbf{1) Базис касательного пространства и нормаль}

\[
\mathbf r_x=(1,0,f_x),\qquad \mathbf r_y=(0,1,f_y).
\]
\[
\mathbf m=\frac{\mathbf r_x\times \mathbf r_y}{\|\mathbf r_x\times \mathbf r_y\|}
=\frac{(-f_x,-f_y,1)}{\sqrt{1+f_x^2+f_y^2}}.
\]

\textbf{2) Касательная плоскость и нормальная прямая}

В точке $P=(x_0,y_0,f(x_0,y_0))$:
\[
z-f(x_0,y_0)=f_x(x_0,y_0)(x-x_0)+f_y(x_0,y_0)(y-y_0).
\]

Нормальная прямая:
\[
(x,y,z)=P+\lambda\bigl(-f_x(x_0,y_0),-f_y(x_0,y_0),1\bigr).
\]

\textbf{3) Первая и вторая фундаментальные формы, гауссова кривизна}

\[
E=1+f_x^2,\qquad F=f_xf_y,\qquad G=1+f_y^2.
\]

Пусть
\[
W=\sqrt{1+f_x^2+f_y^2}.
\]
Тогда
\[
L=\frac{f_{xx}}{W},\qquad M=\frac{f_{xy}}{W},\qquad N=\frac{f_{yy}}{W}.
\]

\[
K=\frac{LN-M^2}{EG-F^2}
=\frac{f_{xx}f_{yy}-f_{xy}^2}{(1+f_x^2+f_y^2)^2}.
\]

\textbf{4) Нормальная кривизна в направлении $a=3\mathbf r_x-2\mathbf r_y$}

\[
k_n(a)=\frac{9L-12M+4N}{9E-12F+4G}
=\frac{9f_{xx}-12f_{xy}+4f_{yy}}{W\bigl(9(1+f_x^2)-12f_xf_y+4(1+f_y^2)\bigr)}.
\]

\textbf{5) Главные кривизны и средняя кривизна}

Главные кривизны являются корнями характеристического уравнения
\[
\det(B-\lambda G)=0,
\]
где
\[
G=\begin{pmatrix}E&F\\F&G\end{pmatrix},\qquad
B=\begin{pmatrix}L&M\\M&N\end{pmatrix}.
\]
Эквивалентно,
\[
k_{1,2}=H\pm\sqrt{H^2-K},
\]
где
\[
H=\frac{EN-2FM+GL}{2(EG-F^2)}
=\frac{(1+f_x^2)f_{yy}-2f_xf_yf_{xy}+(1+f_y^2)f_{xx}}{2(1+f_x^2+f_y^2)^{3/2}}.
\]

\subsection*{e) График функции $z=\sqrt{R^2-x^2-y^2}$ при $x_0=\frac R2,\ y_0=\frac R2,\ z_0=\frac{R\sqrt2}{2}$}

Это верхняя полусфера радиуса $R$. Запишем её как график:
\[
f(x,y)=\sqrt{R^2-x^2-y^2}.
\]
В точке
\[
\left(x_0,y_0\right)=\left(\frac R2,\frac R2\right),\qquad
z_0=f(x_0,y_0)=\frac{R\sqrt2}{2}.
\]

\textbf{1) Базис касательного пространства и нормаль}

\[
f_x=-\frac{x}{\sqrt{R^2-x^2-y^2}},\qquad
f_y=-\frac{y}{\sqrt{R^2-x^2-y^2}}.
\]
В точке:
\[
f_x(x_0,y_0)=f_y(x_0,y_0)=-\frac1{\sqrt2}.
\]

\[
\mathbf r_x=\left(1,0,-\frac1{\sqrt2}\right),\qquad
\mathbf r_y=\left(0,1,-\frac1{\sqrt2}\right),
\]
\[
\mathbf m=\frac{(-f_x,-f_y,1)}{\sqrt{1+f_x^2+f_y^2}}
=\left(\frac12,\frac12,\frac1{\sqrt2}\right).
\]

\textbf{2) Касательная плоскость и нормальная прямая}

\[
\frac12\left(x-\frac R2\right)+\frac12\left(y-\frac R2\right)+\frac1{\sqrt2}\left(z-\frac{R\sqrt2}{2}\right)=0.
\]
После упрощения:
\[
x+y+\sqrt2\,z=2R.
\]

\[
(x,y,z)=\left(\frac R2,\frac R2,\frac{R\sqrt2}{2}\right)
+\lambda\left(\frac12,\frac12,\frac1{\sqrt2}\right).
\]

\textbf{3) Первая и вторая фундаментальные формы, гауссова кривизна}

При этом
\[
E=1+f_x^2=\frac32,\qquad
F=f_xf_y=\frac12,\qquad
G=1+f_y^2=\frac32.
\]

Вторые производные:
\[
f_{xx}=f_{yy}=-\frac{3\sqrt2}{2R},\qquad
f_{xy}=-\frac{\sqrt2}{2R}.
\]
Значит
\[
W=\sqrt{1+f_x^2+f_y^2}=\sqrt2,
\]
\[
L=\frac{f_{xx}}{W}=-\frac{3}{2R},\qquad
M=\frac{f_{xy}}{W}=-\frac{1}{2R},\qquad
N=\frac{f_{yy}}{W}=-\frac{3}{2R}.
\]

\[
K=\frac{LN-M^2}{EG-F^2}=\frac{1}{R^2}.
\]

\textbf{4) Нормальная кривизна в направлении $a=3\mathbf r_x-2\mathbf r_y$}

\[
k_n(a)=\frac{9L-12M+4N}{9E-12F+4G}
=-\frac1R.
\]

\textbf{5) Главные кривизны и средняя кривизна}

\[
k_1=k_2=-\frac1R,\qquad H=-\frac1R,\qquad K=\frac1{R^2}.
\]

\textit{Замечание.} Это та же самая сфера, что и в пункте b), но ориентация нормали здесь задается графиком, поэтому знаки $k_1$, $k_2$ и $H$ получаются противоположными.

\end{document}
